
De @@M0@@, r\'esulte @@M1@@\,; par suite, @@M2@@ n'est pas contenu dans @@M3@@\,; c'est-\`a-dire, la section @@M4@@ n'est pas contenue dans l'image de l'injection @@M5@@, d'o\`u (i).

L'assertion (ii) r\'esulte imm\'ediatement de (i), compte tenu de ce que @@M6@@ implique \'evidemment @@M7@@, car @@M8@@. Dans (iii), on trouve @@M9@@. Donc pour @@M10@@, @@M11@@ car @@M12@@\,; donc @@M13@@ d'apr\`es (ii). Par suite, @@M14@@ pour @@M15@@\,; d'o\`u (iii).

\textbf{Proposition 6.6.} \emph{Sous les conditions de 6.1, on suppose @@M16@@ un @@M17@@-faisceau. Alors @@M18@@ est un @@M19@@-faisceau.}

En effet, soit @@M20@@ une section ``g\'en\'erale'' de @@M21@@ ; alors @@M22@@ d\'efinit un diagramme commutatif \`a lignes et colonnes exactes
@@M23@@
o\`u @@M24@@ est un @@M25@@-faisceau (1.6 (i)) et o\`u @@M26@@ (question de profondeur).
